\chapter{Wu PV to SI Ertel PV: Exact Conversion}
\label{chap:wu_coeffs}

This chapter documents the \emph{exact analytic} conversion from the
Wu Fortran PV scalar $Q_{\text{Wu}}$ to true SI Ertel potential
vorticity.  The conversion coefficient is \textbf{exactly $10^{-8}$}
at every pressure level when $\kappa = 2/7$, which is the standard
dry-air value used throughout the Wu code.  This result was derived on
2026-06-03 and replaces an earlier ad-hoc median-ratio scaling.

\section{Wu PV Definition}

In \texttt{pvpialln\_94UV.f} (lines 460--468), Ertel-like PV is
computed as
\begin{equation}
Q_{\text{Wu}} = -\mathrm{COEF}\cdot\Pi^{-2.5}
  \Bigl[(f+\zeta)\,\frac{\partial\theta}{\partial\Pi}
        + \Bigl(\frac{\partial u}{\partial\Pi}\frac{\partial\theta}{\partial y}
                - \frac{\partial v}{\partial\Pi}\frac{\partial\theta}{\partial x}\Bigr)
  \Bigr],
\label{eq:wu_q}
\end{equation}
where
\begin{align}
\mathrm{COEF} &= \frac{10^{2}\cdot10^{6}\cdot g\cdot\kappa\cdot c_p^{3.5}}{p_0},
\label{eq:coef} \\
\Pi(p) &= c_p\left(\frac{p}{p_0}\right)^{\kappa},
\label{eq:pi_wu}
\end{align}
with numerical constants $g=9.81$\,m\,s$^{-2}$,
$\kappa = 2/7$, $c_p=1004.5$\,J\,kg$^{-1}$\,K$^{-1}$,
$p_0 = 10^{5}$\,Pa (but see warning below).  The ``1.E2'' in the
COEF expression is a legacy over-scaling of no physical significance;
it is the conjugate of the $10^{-2}$ divide applied to PV in the
non-dimensional BALNC/BALP solvers, and it cancels exactly in the
back-conversion derived below.

\section{True SI Ertel PV}

The standard definition of Ertel PV in isobaric coordinates is
\begin{equation}
\mathrm{PV}_{\text{SI}} = -g
  \Bigl[(f+\zeta)\,\frac{\partial\theta}{\partial p}
        + \frac{\partial u}{\partial p}\frac{\partial\theta}{\partial y}
        - \frac{\partial v}{\partial p}\frac{\partial\theta}{\partial x}
  \Bigr],
\label{eq:pv_si}
\end{equation}
with units K\,m$^{2}$\,kg$^{-1}$\,s$^{-1}$
($\equiv 10^{-6}$\,PVU).

\section{Conversion}

The two definitions, \eqref{eq:wu_q} and \eqref{eq:pv_si}, differ only
in the vertical coordinate of differentiation:
$\partial/\partial\Pi$ in Wu versus $\partial/\partial p$ in SI.  The
chain rule relates them:
\begin{equation}
\frac{\partial}{\partial p}
  = \frac{\partial\Pi}{\partial p}\,\frac{\partial}{\partial\Pi}
  = \frac{\kappa\Pi}{p}\,\frac{\partial}{\partial\Pi}.
\label{eq:chain}
\end{equation}
The factor $\kappa\Pi/p$ multiplies \emph{both} the vorticity-stability
term and the shear term identically, so it factors out of the entire
bracket:
\begin{align}
\mathrm{PV}_{\text{SI}}
  &= Q_{\text{Wu}}\cdot\frac{g}{\mathrm{COEF}}\cdot
     \frac{\kappa\Pi}{p}\cdot\Pi^{2.5} \nonumber \\
  &= Q_{\text{Wu}}\cdot\frac{g\,\kappa}{\mathrm{COEF}}\cdot
     \frac{\Pi^{3.5}}{p}
  \equiv Q_{\text{Wu}}\cdot c(p).
\label{eq:exact}
\end{align}

\subsection{The Coefficient $c(p)$ is Constant}

Substituting $\Pi = c_p(p/p_0)^{\kappa}$ and
$\mathrm{COEF} = 10^{8}\,g\,\kappa\,c_p^{3.5}/p_0$:
\begin{align}
c(p) &= \frac{g\,\kappa}{10^{8}\,g\,\kappa\,c_p^{3.5}/p_0}\cdot
       \frac{c_p^{3.5}(p/p_0)^{3.5\kappa}}{p} \nonumber \\
     &= \frac{1}{10^{8}}\,(p/p_0)^{3.5\kappa-1}.
\label{eq:c_p}
\end{align}
With $\kappa = 2/7$,
\begin{equation}
3.5\kappa = \frac{7}{2}\cdot\frac{2}{7} = 1
\quad\Longrightarrow\quad
c(p) = \frac{1}{10^{8}}\,(p/p_0)^{0} = \mathbf{10^{-8}}
\label{eq:constant}
\end{equation}
at \emph{every} pressure level.  Wu's choice of the $\Pi^{-2.5}$
weighting was designed precisely so that $\Pi^{3.5}/p$ is
pressure-independent: the $p$ dependence in $\Pi^{3.5}$ exactly
cancels with the $1/p$ from the chain rule.

\begin{keybox}[Wu Q $\to$ SI conversion is a universal constant]
  \[
  \boxed{\;\mathrm{PV}_{\text{SI}} = Q_{\text{Wu}} \times 10^{-8}\;}
  \qquad\text{or}\qquad
  \boxed{\;Q_{\text{Wu}} = \mathrm{PV}_{\text{SI}} \times 10^{8}\;}
  \]
  This holds for \emph{any} pressure level when $\kappa = 2/7$.
  If a different $\kappa$ were used, the per-level array
  $c(p) = (g\kappa/\mathrm{COEF})\cdot\Pi^{3.5}/p$ must be
  evaluated; the `WU2SI` array in \texttt{parse\_and\_pv.py}
  implements this general form.
\end{keybox}

\section{Agreement with ERA5 (after the raw-$T$ fix)}

Dividing Wu $Q$ by $10^{8}$ converts it to \emph{genuine SI units}
(K\,m$^{2}$\,kg$^{-1}$\,s$^{-1}$).  Once the input bug was fixed --- the
grid writer now feeds \texttt{pvpialln} \textbf{raw temperature $T$}, not
$\theta$, so the Exner conversion is applied exactly once --- the
converted field agrees with ERA5's native Ertel PV essentially one to
one.  At 250\,hPa the current pipeline reports
\begin{equation}
\text{median ratio } \frac{Q_{\text{Wu}}\times10^{-8}}{\text{PV}_{\text{ERA5}}} = 1.00,
\qquad \text{slope} = 0.99,\qquad \corr = 0.995 .
\end{equation}

\begin{warnbox}[History: the spurious 3.5x residual]
An earlier version fed \emph{$\theta$} into \texttt{pvpialln}, which then
applied the Exner factor a \emph{second} time.  This double conversion
inflated the static stability $\partial\theta/\partial\Pi$ and hence $Q$
by a level-dependent $\sim$2--4$\times$ (median $\sim$3.5$\times$), and
was masked at the time by an ad-hoc \texttt{STAB\_SCALE}$\approx$0.3.
Both the double conversion and the fudge factor are gone; the residual
factor is now $\approx 1$.  No empirical scaling is applied --- only the
analytic $10^{-8}$.
\end{warnbox}

\noindent (Equivalently, since $\mathrm{PVU}=10^{-6}$\,SI, the Wu scalar
is $\approx 100\times$ a PVU value: $\mathrm{PVU}=Q_{\text{Wu}}\times10^{-2}$.)

\section{Implementation in the Pipeline}

In \texttt{wu/steps/08\_parse\_outputs/parse\_and\_pv.py}:
\begin{lstlisting}[language=Python]
# Exact Wu Q -> SI Ertel PV per-level coefficient
PLEVS_PA = np.asarray(WU_LEVS) * 100.0
PI_LEV   = CP * (PLEVS_PA / P0) ** KAP
COEF_WU  = 1.0e2 * 1.0e6 * 9.81 * KAP * CP**3.5 / P0
WU2SI    = (G * KAP / COEF_WU) * PI_LEV**3.5 / PLEVS_PA  # ~1e-8
Q_wu_si  = Q_wu * WU2SI[k250]   # SI [K.m2.kg-1.s-1]
Q_wu_pvu = Q_wu_si * 1e6        # PVU for display
\end{lstlisting}

The $c(p) = 10^{-8}$ constant is also documented in
\texttt{config.py} section~7 alongside the full Wu non-dimensional
scale set (DPI, LL, FF, THO, HND, QCONST, PI\_WU, PIF, BB/BH/BL).

\begin{warnbox}[COEF uses Fortran g=9.81]
The Wu Fortran hardcodes $g=9.81$ (line 460), while the standard
SI value is $g=9.80665$\,m\,s$^{-2}$.  The Python coefficients
match the Fortran's 9.81 for exact numeric consistency.  The
difference is $<0.04\%$ and is absorbed by the residual
discretisation bias anyway.
\end{warnbox}
